% !TeX root = ../papers/nc-p3.tex
\section{平行二重色の残部に対する四次・七次障害}
\label{sec:parallel-excluded-shapes}
\begin{theorem}
平行二重色の代表台のうち $67$--$75$, $86,87$, $91$--$93$, $97$--$111$ の29台には標準AS窓がない。
さらに
\[
 79,83,123,128,130,131,137,141
\]
の八台にもなく、
$80,81,84,119,120,121,122,125,126,127,132,133,135,136,138,140$ の16台には指定階数の係数すらない。
従って定理\ref{thm:klein-parallel-flags}, \ref{thm:parallel-symmetric-surface}と合わせ、
148台のうち、この段階で正則性に未決定部分が残るのは
\[
 82,85,88,89,90,118,124,129,134,139,142,143,144,145,146,147
\]
だけである。
\end{theorem}
\begin{proof}
定理\ref{thm:klein-parallel-flags}のSmith解で、最後の16台では残る小行列式の少なくとも一つが各位相で
Laurent単元となる。29台では同じ単元障害と、自由トーラス上での単元二列小行列式を使った四次障害に加え、
残る係数方程式の解での検査が必要となる。全ての位相を保持して以下の操作を行う。

座標は全て可逆である。可逆と分かった因子を除いた各方程式を既約因子の積とし、ある因子集合が別のものを
含むなら大きい方の方程式を落とす。複数因子の式では最初に零となる因子で分ける。一因子の式 $Ay+B=0$ があれば、
$A=B=0$ の部分と $AB\ne0$, $y=-B/A$ の部分に分ける。以前に除いた因子と代入の分母も追跡する。
よって低次元の特殊化も被覆する。検算は各段での因子分解と、末端で全ての元の式が恒等的に消えることを確かめる。
107番だけは五座標を順に $a,b,c,d,e$ と書き、
\[
 c=d(ab+1)/a,\qquad e^2=(ab+1)(bd+1)/a^2
\]
を先に解く。残る条件は $(a-1)(ab+1)d+a^2=0$ であり、全係数非零だから
$d=-a^2/((a-1)(ab+1))$, $b=(a-1)(1-a^2e^2)/a$ の一つの有理表示で全て覆える。

各末端で実際の六関係の列空間を有理関数体上で取る。始点を固定した零文字の長さ4の関係挿入行列の
余次元が9を超えれば、分母が定義されたその表示のどの点でも同じ下界を持つ。
実際、全ての列を使う元の語基底での挿入行列の小行列式が恒等的に消えることを、関数体で基底を選んで調べただけである。
正則な窓の期待文字次元はKoszul分解から、次数4から7まで順に
$(9,10,8,8),(14,14,14,14),(20,20,22,22),(30,30,30,30)$ となる。
29台は階数不足または四次障害までで全て排除される。

残り八台には同じ全係数被覆の上で、直前の経路基底に一文字を付け、末尾に二次関係を挿入することで
有理関数体上の高次商を作る。各段の文字別の自由列を新しい経路基底とし、全体を分数なしの消去で計算する。
$79,123,128$ は始点0・長さ7・文字0、$83,137,141$ は始点3・長さ7・文字1、
$130$ は始点1・長さ5・文字1、$131$ は始点2・長さ6・文字3で期待値を超える。ここでも元の語への
関係挿入は係数について正則なので、全ての特殊化まで排除される。
\path{computations/verify_parallel_excluded_shapes.py} が全位相と係数分岐に対してこれらの不等式を確認する。
\end{proof}

\begin{proposition}
定理\ref{thm:parallel-symmetric-surface}の二曲面と定理\ref{thm:klein-three-lines}の二族との間に無印の重なりはない。
\end{proposition}
\begin{proof}
反対階数が一致する候補は対称表示同士だけである。三薄色の対称表示では、
定理\ref{thm:three-lines-symmetric-isomorphism}のFourier座標において二次式の束は跡零対角網と第四行列の和になる。
網の射影平面には六つの階数2点があり、第四行列の跡は非零である。
網を跡零の形に規格化すると、その安定群は共通倍率を除き符号付き置換であるため、この跡の零・非零は網を指定した合同の不変量である。

平行曲面の行列 $H$ を前節と同じ記号で書く。階数2点は二平面 $s=0$ と $t=0$ の和に入る。
$s=0$ には $YT$ と、$a=0$ または $b=0$ に各高々二点、さらに $pq=1$ のときだけ $a=1,b=-p,t=0$ がある。
$pq=1$ ならこの平面での行列式は $ab((pa+b)^2/p-t^2)$ で、四直線の相異なる六交点となる。
この場合を含め平面内のどの直線にも三点以下である。$t=0$ のうち $s\ne0$ の点は、
$a=1$, $b^2+(p-q)b-1=0$ の高々二通りの $b$ で、各々一本の直線上に二点以下である。
この二直線は交わり、他の直線は各一本と一点でしか交わらない。従って上の二平面以外の平面には合計五点以下しかない。
$t=0$ で $s=0$ の座標軸点が二つ同時に加わるのは $p=q=0$ のときだけであり、
$pq=1$ の場合 $b=-p$ の根の $s$ は零なので $t=0$ の点は計三点以下である。

よって六点平面の候補は $pq=1$ の $s=0$ と $p=q=0$ の $t=0$ だけである。
前者は $X,\sqrt p Z,\sqrt p Y+T,\sqrt p Y-T$ の座標で跡零対角網になり、第四式 $XT-YZ$ は異なる二組をまたぐため跡零である。
後者の網は $(X,T)$ と $(Z,Y)$ の二ブロックに分かれ、行列式は $(a^2+s^2)(b^2+s^2)$。
四直線に沿って各ブロックを同時対角化し、各座標を倍率変更すれば跡零対角網になる
（四対角値の唯一の線型関係には全四係数が現れる）。第四式 $YT$ は再び組をまたいで跡零である。
従っていずれも三薄色の網には合同にならない。
\end{proof}
